% chapters/01-introduction.tex
\chapter{บทนำ}
\label{ch:intro}

\section{ความเป็นมาและความสำคัญของปัญหา}
\label{sec:intro-background}

แพลตฟอร์มวิดีโอออนไลน์เป็นสื่อหลักของผู้บริโภคในยุคปัจจุบัน ในปี ค.ศ.\ 2025
YouTube มีผู้ใช้งานต่อเดือนมากกว่า 2.5 พันล้านคนทั่วโลก \cite{kemp2025digital}
และจัดเป็นแพลตฟอร์มวิดีโอที่มีการเข้าถึงสูงสุดในประเทศไทย รายงาน Digital
Thailand 2025 ระบุว่าผู้ใช้อินเทอร์เน็ตชาวไทยรับชม YouTube เฉลี่ย 1 ชั่วโมง
44 นาทีต่อวัน \cite{datareportal2025thailand} การเติบโตของระบบนิเวศการสร้าง
รายได้ผ่าน YouTube Partner Programme รวมถึงการรับสปอนเซอร์โดยตรง การตลาดผ่าน
ผู้มีอิทธิพล (influencer marketing) และการขายสินค้าผ่านวิดีโอ ทำให้คุณค่าทาง
เศรษฐกิจของแต่ละวิดีโอที่ \emph{เป็นกระแส} หรือ \emph{ไวรัล} (viral) มีนัยสำคัญ
อย่างยิ่งทั้งต่อผู้สร้างเนื้อหารายเล็กและบริษัทขนาดใหญ่

ในขณะเดียวกัน อัตราการแข่งขันบนแพลตฟอร์มเพิ่มสูงขึ้นเรื่อย ๆ มีการประเมินว่ามี
วิดีโอใหม่ถูกอัปโหลดขึ้น YouTube กว่า 500 ชั่วโมงต่อนาที \cite{ceci2024youtube}
สัดส่วนของวิดีโอที่กลายเป็นกระแสในวงกว้างจึงเหลือเพียงเศษเสี้ยวของวิดีโอทั้งหมด
ผู้สร้างเนื้อหา ผู้จัดการช่อง และฝ่ายการตลาดในประเทศไทยยังต้องอาศัย
\emph{การคาดเดาเชิงประสบการณ์} เป็นหลักในการตัดสินใจเลือกหัวข้อ พาดหัว ปก และ
เวลาเผยแพร่ ซึ่งทำให้การลงทุนในแต่ละวิดีโอมีความเสี่ยงสูงและขาดข้อมูลเชิงประจักษ์
มาสนับสนุน

ในมุมมองทางวิทยาศาสตร์ข้อมูล การพยากรณ์ความไวรัลของวิดีโอ \emph{ก่อน} เผยแพร่
เป็นปัญหาที่ท้าทาย เพราะข้อมูลที่ใช้พยากรณ์ในขณะเผยแพร่มีจำกัดอยู่ที่
\emph{ข้อความ} เป็นหลัก ได้แก่ ชื่อวิดีโอ คำอธิบาย แท็ก หมวดหมู่ และข้อมูล
\emph{เชิงโครงสร้าง} ของช่อง (เช่น จำนวนผู้ติดตาม วันที่สร้างช่อง อัตราการ
อัปโหลดเฉลี่ย) เนื้อหาที่อยู่ในรูปข้อความภาษาไทยมีลักษณะเฉพาะที่ต่างจากภาษา
ตะวันตก เช่น ไม่มีเครื่องหมายเว้นวรรคระหว่างคำ การปะปนกับภาษาอังกฤษและภาษา
ในกลุ่มประเทศเพื่อนบ้าน การใช้อิโมจิและเครื่องหมายพิเศษเพื่อดึงความสนใจ ทำให้
การใช้แบบจำลองภาษา (Language Models, LMs) ที่ผ่านการฝึกล่วงหน้าด้วยข้อมูล
ภาษาไทยขนาดใหญ่จึงเป็นทางเลือกที่น่าสนใจ

\emph{แบบจำลองทรานส์ฟอร์เมอร์} (Transformer) ที่นำเสนอโดย \textcite{vaswani2017attention} ได้กลายเป็นแกนสำคัญของระบบ NLP สมัยใหม่ เนื่องจาก
สถาปัตยกรรม self-attention แก้ปัญหา \emph{การเลือนหายของเกรเดียนต์} (vanishing
gradient) ที่ RNN/LSTM ประสบเมื่อต้องจดจำบริบทระยะยาว และสามารถจับการพึ่งพา
ระยะไกลในข้อความ (long-range dependencies) ได้อย่างมีประสิทธิภาพ \textcite{devlin2019bert} เสนอ BERT ซึ่งเป็นแบบจำลองแบบเข้ารหัส (encoder) สองทิศทาง
ที่ฝึกล่วงหน้าด้วยภาระงาน Masked Language Modelling (MLM) และต่อมา \textcite{liu2019roberta} ได้พัฒนา RoBERTa ซึ่งปรับขั้นตอนการฝึกล่วงหน้าให้แข็งแกร่ง
ขึ้น แบบจำลองในตระกูล BERT/RoBERTa สำหรับภาษาไทยที่ได้รับความนิยมในปัจจุบัน
ได้แก่ \textbf{WangchanBERTa} โดย AIResearch.in.th \cite{lowphansirikul2021wangchanberta}
และ \textbf{PhayaThaiBERT} โดยทีม \textcite{phayathaibert2023} ส่วน
แบบจำลองมัลติลิงกัวลที่รองรับภาษาไทยและสามารถจัดการกับการสลับภาษาในเอกสารเดียว
กันได้ดี คือ \textbf{XLM-RoBERTa} \cite{conneau2020xlmr}

อย่างไรก็ตาม \emph{ยังไม่มีงานวิจัยใดในประเทศไทย} ที่เปรียบเทียบความสามารถใน
การพยากรณ์ความไวรัลของวิดีโอ YouTube ระหว่างแบบจำลองทรานส์ฟอร์เมอร์ภาษาไทย
หลายตัวบนชุดข้อมูลเดียวกัน พร้อมการประเมินผลทางสถิติที่เข้มงวด งานก่อนหน้า
ในต่างประเทศ เช่น \textcite{trzcinski2017popularity} และ
\textcite{khosla2014popularity} ใช้คุณลักษณะภาพเป็นหลัก ในขณะที่
ไม่ครอบคลุมลักษณะเฉพาะของข้อความภาษาไทย ดังนั้นการศึกษานี้จึงเข้าไปอุดช่องว่าง
ดังกล่าวด้วยการเปรียบเทียบที่ครอบคลุมและมีระเบียบวิธีวิจัยที่สามารถทำซ้ำได้

\section{วัตถุประสงค์ของการวิจัย}
\label{sec:intro-objectives}

\begin{enumerate}[leftmargin=2.5em,label=\arabic*.]
  \item เพื่อพัฒนากระบวนการรวบรวมและจัดเตรียมข้อมูลวิดีโอภาษาไทยจาก YouTube
        Data API v3 ที่สามารถทำซ้ำได้ และครอบคลุมข้อมูล \emph{เชิงข้อความ}
        และ \emph{เชิงโครงสร้าง} ของวิดีโอและช่อง
  \item เพื่อกำหนด \emph{ดัชนีความไวรัล} (Virality Index) ระดับช่องที่
        ปราศจากความเอนเอียงเชิงขนาดช่อง และแปลงเป็นปัญหาการจำแนกประเภทไบนารี
        แบบ top-decile per channel
  \item เพื่อเปรียบเทียบประสิทธิภาพของแบบจำลองทรานส์ฟอร์เมอร์ภาษาไทย 3 แบบ
        ได้แก่ WangchanBERTa, PhayaThaiBERT และ XLM-RoBERTa-large บน
        ปัญหาเดียวกัน ด้วย ROC-AUC, F1, PR-AUC, MCC พร้อมช่วงความเชื่อมั่น
        แบบบูตสแตรป
  \item เพื่อเปรียบเทียบแบบจำลองทรานส์ฟอร์เมอร์กับแบบจำลองพื้นฐานเชิงคลาสสิก
        (Logistic Regression, LightGBM, XGBoost) และแบบจำลองไฮบริด (MLP +
        LightGBM head) ที่ใช้คุณลักษณะร่วม (combined features)
  \item เพื่อพัฒนาและประเมิน \emph{การรวมแบบจำลอง} (stacking ensemble) ที่รวม
        แบบจำลองฐาน 15 ตัว และเปรียบเทียบประสิทธิภาพหลังจาก
        \emph{การปรับเทียบความน่าจะเป็น} (probability calibration) ด้วย
        Platt scaling และ Isotonic regression
  \item เพื่อทดสอบทางสถิติว่ารูปแบบความผิดพลาดของแบบจำลองทั้งหลายมีความแตกต่าง
        อย่างมีนัยสำคัญหรือไม่ ด้วย McNemar's test แบบคู่ (pairwise) และ
        Cochran's Q test ข้ามแบบจำลองหลายตัวบนชุดทดสอบเดียวกัน
  \item เพื่อวิเคราะห์ \emph{ความสามารถในการอธิบาย} (explainability) ของ
        แบบจำลองด้วย SHAP, LIME และ Attention Rollout
\end{enumerate}

\section{ขอบเขตของการวิจัย}
\label{sec:intro-scope}

\subsection{ขอบเขตด้านปัญหา}
\begin{enumerate}[leftmargin=2.5em,label=\arabic*.]
  \item ปัญหาที่ศึกษาเป็น \textbf{การจำแนกประเภทไบนารี} (binary classification)
        ของวิดีโอภาษาไทยว่าจะเป็นไวรัลหรือไม่ ไม่ใช่ปัญหาการถดถอย (regression)
        ของจำนวนยอดเข้าชม
  \item ใช้ข้อมูลที่สามารถสังเกตได้ \emph{ในขณะเผยแพร่} เป็นหลัก ได้แก่
        ชื่อวิดีโอ คำอธิบาย แท็ก หมวดหมู่ ฟีเจอร์ของช่อง และเวลาเผยแพร่
        ไม่รวมข้อมูลภายหลังเผยแพร่ที่ผู้ใช้ทั่วไปไม่สามารถเห็น เช่น watch
        time, audience retention curve
  \item ดัชนีความไวรัลคำนวณจากตัวชี้วัดการมีส่วนร่วม 30 วันหลังเผยแพร่ ไม่ใช่
        ตัวชี้วัดสุดท้ายของช่วงอายุของวิดีโอ
\end{enumerate}

\subsection{ขอบเขตด้านแบบจำลอง}

งานวิจัยนี้กำหนด \emph{ขอบเขตการเปรียบเทียบหลัก} ที่แบบจำลองทรานส์ฟอร์เมอร์
แบบ \textbf{เข้ารหัส} (encoder) เนื่องจากปัญหาเป็นการจำแนกประเภท ไม่ใช่งาน
สร้างข้อความ (generation) แบบจำลอง 3 ตัวที่เปรียบเทียบ ได้แก่
WangchanBERTa, PhayaThaiBERT และ XLM-RoBERTa-large

หมายเหตุ: แบบจำลองแบบ \emph{ถอดรหัส} (decoder-only) เช่น Typhoon 2.5 และ
OpenThaiGPT ซึ่งเคยอยู่ในแผนเริ่มต้น ถูกตัดออกจากการเปรียบเทียบหลักด้วยเหตุผล
3 ประการคือ (1) ผลการศึกษาเชิงทฤษฎีและเชิงประจักษ์โดย \textcite{touvron2023llama} และ \textcite{wei2022emergent} แสดงว่า decoder
LLM เหมาะกับงาน generation มากกว่างาน classification ขนาดเล็ก (2) ทรัพยากร
คอมพิวเตอร์ที่จำกัดทำให้การปรับจูนเต็มรูปแบบเป็นไปไม่ได้ ต้องใช้ QLoRA ซึ่ง
เป็นการเปรียบเทียบที่ไม่ยุติธรรมเชิงสถาปัตยกรรม และ (3) ข้อสังเกตเชิงประจักษ์
ในชุมชน NLP ระบุว่า encoder ขนาดกลางยังเป็นทางเลือกที่ดีกว่าในเชิงต้นทุน--
ประสิทธิผลสำหรับงาน text classification ขนาดเล็ก
อย่างไรก็ตามรหัสฐาน (\texttt{src/models/transformer\_finetune.py}) ยังคงรองรับ
เส้นทาง QLoRA สำหรับ decoder LLM เพื่อการขยายผลงานในอนาคต

\subsection{ขอบเขตด้านข้อมูล}
\begin{itemize}[leftmargin=2em]
  \item ภาษา: ไทย (รวมถึงเอกสารที่ปะปนกับภาษาอังกฤษและอีโมจิ)
  \item ช่อง: 82 ช่อง คัดเลือกตามหมวดหมู่หลัก 3 ประเภท (Entertainment, Gaming,
        Music) และขนาดช่อง 3 ระดับ (mid 100k--1M, large 1M--10M, mega $\geq$10M)
  \item ขนาด: 23{,}431 วิดีโอ
  \item ช่วงเวลา: เผยแพร่ระหว่าง พ.ศ.\ 2562--2569 (10 ปี)
  \item อัตราส่วนคลาสบวก: 10.16\% ตามนิยาม top-decile per channel
\end{itemize}

\section{ประโยชน์ที่คาดว่าจะได้รับ}
\label{sec:intro-benefits}

\begin{enumerate}[leftmargin=2.5em,label=\arabic*.]
  \item ผู้สร้างเนื้อหาภาษาไทยและฝ่ายการตลาดได้รับเครื่องมือพยากรณ์ที่ผ่านการ
        ตรวจสอบทางสถิติ สามารถนำไปทดสอบหัวข้อและพาดหัวก่อนเผยแพร่จริง
  \item ชุมชน NLP ภาษาไทยได้รับ \emph{เกณฑ์มาตรฐาน} (benchmark) ที่เปรียบเทียบ
        ทรานส์ฟอร์เมอร์ภาษาไทย 3 ตัวบนปัญหาการพยากรณ์ความไวรัล พร้อมการ
        เปิดเผยการทำนายในรูป Parquet เพื่อให้นักวิจัยรายอื่นเปรียบเทียบได้
  \item ผู้วิจัยและนักศึกษาด้านสถิติประยุกต์ได้รับ \emph{กรอบการวิเคราะห์เชิง
        สถิติเข้มงวด} (McNemar / Cochran's Q / bootstrap CI / calibration)
        ที่สามารถนำไปประยุกต์กับปัญหาการจำแนกประเภทอื่น
  \item รหัสฐาน (codebase) เป็น open-source ที่ทำซ้ำได้ตั้งแต่ \texttt{make
        prepare} จนถึง \texttt{make eval} โดยมี MLflow และ git SHA tagging
        ทุกการรัน
\end{enumerate}

\section{นิยามศัพท์ที่ใช้ในงานวิจัย}
\label{sec:intro-terms}

\begin{description}[leftmargin=2.5em,style=nextline]
  \item[ความไวรัล (Virality)]
    คุณสมบัติของวิดีโอที่มีการแพร่กระจายและการมีส่วนร่วมของผู้ชมในระดับสูง
    เกินกว่าค่ามาตรฐานของช่อง ไม่ได้หมายถึงจำนวนยอดเข้าชมที่สูงเพียงอย่างเดียว
    เนื่องจากช่องขนาดใหญ่มียอดเข้าชมพื้นฐานสูงอยู่แล้ว ดังนั้นต้องนิยามให้
    \emph{สัมพัทธ์ภายในช่องเจ้าของ}
  \item[ดัชนีความไวรัล (Virality Index, $\mathit{VI}$)]
    คะแนนต่อเนื่องที่กำหนดให้แต่ละวิดีโอ คำนวณจากผลรวมแบบถ่วงน้ำหนักของอัตราส่วน
    การมีส่วนร่วม อัตราส่วนการกดถูกใจ และอัตราส่วนความคิดเห็น หลังการแปลง
    ลอการิทึมและ z-score ภายในช่อง (รายละเอียดในบทที่ \ref{ch:method})
  \item[ป้ายกำกับไวรัล (Viral Label)]
    ตัวบ่งชี้ไบนารีที่กำหนดเป็น 1 หากวิดีโอมี $\mathit{VI}$ อยู่ใน 10\% บนสุด
    ของช่องเจ้าของ และ 0 ในทางตรงข้าม
  \item[การปรับจูน (Fine-tuning)]
    การฝึกแบบจำลองทรานส์ฟอร์เมอร์ที่ผ่านการฝึกล่วงหน้าแล้ว ต่อด้วยข้อมูลของ
    งานปลายน้ำ (downstream task) เพื่อปรับน้ำหนักให้เข้ากับปัญหาเฉพาะ
  \item[การรวมแบบจำลอง (Stacking Ensemble)]
    เทคนิคที่ใช้ \emph{แบบจำลองระดับเมตา} (meta-learner) ในการรวมความน่าจะเป็น
    จากแบบจำลองฐานหลายตัว เพื่อให้ได้พยากรณ์รวมที่ดีกว่าเดี่ยว ๆ
    \cite{wolpert1992stacked}
  \item[การปรับเทียบความน่าจะเป็น (Probability Calibration)]
    กระบวนการแปลงคะแนนพยากรณ์ของแบบจำลองให้สอดคล้องกับความน่าจะเป็นจริงของ
    เหตุการณ์ วิธีที่ใช้กันแพร่หลายคือ Platt scaling \cite{platt1999platt}
    และ Isotonic regression \cite{zadrozny2002isotonic}
\end{description}

\section{สมมติฐานของการวิจัย}
\label{sec:intro-hypotheses}

จากการสำรวจวรรณกรรมในบทที่ \ref{ch:litreview} คณะผู้วิจัยตั้งสมมติฐานเชิง
ปฏิบัติการ (operational hypotheses) ดังนี้

\begin{description}[leftmargin=2.5em,style=nextline]
  \item[$H_1$:] PhayaThaiBERT ซึ่งมีพารามิเตอร์ 278M และฝึกบนคลังภาษาไทยที่ใหญ่
    กว่า จะให้ ROC-AUC สูงกว่า WangchanBERTa ซึ่งมี 105M
  \item[$H_2$:] XLM-RoBERTa-large ซึ่งมีพารามิเตอร์ 560M และเป็นแบบจำลอง
    มัลติลิงกัวล จะให้ประสิทธิภาพไม่แย่กว่าแบบจำลอง monolingual ภาษาไทย เนื่องจาก
    การสนับสนุนการสลับภาษา (code-switching) ในชื่อวิดีโอ YouTube ภาษาไทย
  \item[$H_3$:] การรวมแบบจำลองระหว่างแบบจำลองเชิงคลาสสิกและทรานส์ฟอร์เมอร์
    จะให้ ROC-AUC สูงกว่าแบบจำลองเดี่ยวที่ดีที่สุดอย่างมีนัยสำคัญทางสถิติ
    ($p < 0.05$ จาก McNemar's test)
  \item[$H_4$:] รูปแบบความผิดพลาดของทรานส์ฟอร์เมอร์ทั้งสามมีความแตกต่างกันอย่าง
    มีนัยสำคัญ (Cochran's Q $p < 0.05$) สนับสนุนการใช้การรวมแบบจำลองหลาย
    สถาปัตยกรรม
\end{description}

ผลการตอบสมมติฐานจะรายงานในบทที่ \ref{ch:results}
